For $p=2$, differentiate the symbol and use spatial Plancherel. Changing variables to $u=|\xi|^2/r^2$ in the inner integral gives
\[
\|\mathcal G_\alpha f\|_2^2=
2\alpha^2\int_0^1u(1-u)^{2\alpha-2}\,du\,\|f\|_2^2
=\frac{\alpha}{2\alpha-1}\|f\|_2^2.
\]
Suppose $p>2$. Set $s_+(p)=\max\{0,(n-1)/2-n/p\}$ and choose $\varepsilon>0$ so that $\sigma=s_+(p)+\varepsilon<\alpha-1/2$. On $r=2^k s$, $1\le s\le2$, put $\rho=2^{-k}|\xi|$ and use $dr/r=ds/s$, a measure comparable to $ds$. The multiplier of $r\partial_r B_r^\alpha$ is zero for $\rho>s$ and equals
\[
 m(s,\rho)=2\alpha \rho^2s^{-2\alpha}(s+\rho)^{\alpha-1}(s-\rho)^{\alpha-1}
 \qquad (\rho<s).
\]
This follows by differentiating the locally integrable positive-part power, which has no boundary delta term since $\alpha>0$.

Take a smooth cutoff equal to one near $\rho=s$ and supported where $|s-\rho|<1/2$ for $s\in[1,2]$. Insert a smooth cutoff in $s$ agreeing with one on $[1,2]$, with support in $(3/4,9/4)$. The factors preceding $(s-\rho)^{\alpha-1}$ times the cutoffs give a smooth compactly supported amplitude $A(s,\rho)$ with $\rho$ bounded away from zero. For $s-\rho>0$ split this boundary term using a smooth dyadic partition $\sum_j\eta(2^j(s-\rho))=1$ with $\eta\in C_c^\infty((1/2,2))$. The support leaves only $j\ge0$. With $\delta=2^{-j}$ and $\varphi(u)=\eta(u)u^{\alpha-1}$ on its support, each piece has symbol
\[
 \delta^{\alpha-1} A(s,\rho)\varphi((s-\rho)/\delta).
\]

For each $\delta$, let $T_\delta^{(k)}$ be the spatial dilation of the operator in \ref{ms:lemma-11-2}, with $s\in[1,2]$ unchanged. By integration by parts up to $N=n+3$ times in $\xi$ on a fixed compact annulus, its undilated kernel $K$ and first spatial derivatives have decay $C\delta^{-N}(1+|x|)^{-N}$ uniformly in $s$. In particular the same estimate holds for their norms in $H=L^2([1,2],ds)$ and $H^n$. Choose smooth annular projections $P_k$, dilates of one multiplier equal to one on the undilated frequency support. Apply \ref{ms:lemma-11-3} to the inputs $P_kf$, using the bounds in \ref{ms:lemma-11-2} for $B_p$ and $B_2$. Proposition~\ref{ext:stein-annular-square} bounds the input square norm by $C\|f\|_p$ for the chosen smooth annular projections. On all dyadic intervals of radii together the shell contributes at most
\[
 C\delta^{\alpha-1}
 \left(\delta^{1/2}+\delta+
 [1+\log(1/\delta)]^{1/2-1/p}\delta^{1/2-\sigma}\right)\|f\|_p.
\]
The sum over all $j$ converges since $\alpha-1/2-\sigma>0$.

The multiplier remaining away from the boundary is smooth, supported in $\rho\le2$ for $s\in[1,2]$, and vanishes to order two at zero. Split it by smooth annular cutoffs at $\rho\asymp2^\ell$, $\ell\le2$. Dividing each term by $2^{2\ell}$ and rescaling $\rho$ to the unit annulus gives symbols with uniformly bounded derivatives of every fixed order. The undilated kernels and their first derivatives then satisfy \ref{ms:lemma-11-3} at $\delta=1$ with uniform $C_*$. The same kernel decay and convolution bound give uniform $B_2$ and $B_p$. For each $\ell$ apply that lemma using the projections at scales $2^{k+\ell}$. Proposition~\ref{ext:stein-annular-square} with shift $\ell$ bounds their input square norm uniformly in $\ell$. The resulting bounds sum because $\sum_{\ell\le2}2^{2\ell}<\infty$.

These estimates apply initially to finite sets of $k$ and finite shell sums. For Schwartz $f$ the multiplier integrals converge to the original integral for each fixed $r$ by local integrability of the power $(s-\rho)^{\alpha-1}$ and bounded overlap in the partitions. Fatou's lemma in the square integral and then in $x$ gives the bound on any finite family of radius intervals. Monotone convergence gives the full estimate. Density extends the linear map $f\mapsto (r\partial_r B_r^\alpha f)_{r>0}$ to a bounded map into $L_x^p(L^2(dr/r))$ and therefore extends its pointwise Hilbert norm.
